
\section{Blocks of maximal defect: reduction to \texorpdfstring{$G/O_{2'}(G)$}{G/O2'(G)}}\label{sec:4}%



\noindent First, we use Fong's reductions to reduce the classification of blocks of maximal defect with dihedral defect groups up to splendid Morita equivalence to cyclic $p'$-central extensions of groups $G$ with dihedral Sylow $2$-subgroups and $O_{2^\prime}(G)=1$. 

\begin{lem}\label{lem:Fong2ndredDihedral}%
     Let $G$ be a finite group with dihedral Sylow $2$-subgroups of order $2^n$, where $n\geq 3$. Then any block $\bB\in\Bl^{\max}(kG)$ is splendidly Morita equivalent to a block $\widehat{\bB}\in\Bl_{\max}(k\widehat{L})$, where $L$ is a finite group with  $O_{2'}(L)=1$ and dihedral Sylow $2$-subgroups of order $2^n$,  and $\widehat{L}$ is a central extension of $L$ by a cyclic group of order $|L|_{2'}$. Moreover, if $H^2(L,k^\times)=1$, then $\widehat{L}$ may be replaced by $L$ itself.
\end{lem}


\begin{proof}
We argue by induction on $|G|$. If $|G|=8$, then $G\cong D_8$ and $\bB=kG$.
Taking $L=\widehat{L}=G$ and $\widehat{\bB}=\bB$ proves
the assertion in this case.
Now assume that $|G|>8$ and that the assertion holds for every group satisfying the hypotheses whose order
is strictly smaller than $|G|$.
Set $N:=O_{2'}(G)$,  let $D\in\Syl_2(G)$, let $\bB\in\Bl^{\max}(kG)$. Let $\bc\in\Bl(kN)$ be a block covered by $\bB$  and let $T$ denote its stabilizer in $G$. 
Suppose first that $T<G$.
By the Fong{--}Reynolds Theorem (Theorem~\ref{thm:FongReynolds}) the block $\bB$
is splendidly Morita equivalent to a block $\widetilde{\bB}\in\Bl(kT)$. Moreover, $\widetilde{\bB}$ has a defect group $D'\leq T$
which is $G$-conjugate to $D$.
Since $\bB$ has maximal defect, $D'$ is a dihedral Sylow $2$-subgroup of $T$
of order $2^n$, and $\widetilde{\bB}$ has maximal defect.
The assertion follows by applying the induction
hypothesis to $T$ and $\widetilde{\bB}$, and using the
transitivity of splendid Morita equivalences.

We may therefore assume that $T=G$, so that $\bc$ is
$G$-stable.
Set $L:=G/N$. Then $O_{2'}(L)=1$. Since $D\cap N=1$, the subgroup $DN/N\in\Syl_2(L)$
is isomorphic to $D$. 
 Then, Theorem~\ref{thm:Fong2ndReduction} (Fong's 2nd reduction) yields 
existence of a Morita equivalence between $\bB$ and
a block $\widehat{\bB}\in\Bl_{\max}(k\widehat{L})$,
where $\widehat{L}$ is a central extension of $L$
by a cyclic group of order $|L|_{2'}$.
Moreover, after identifying the defect groups, this
Morita equivalence is induced by a bimodule with
vertex $\Delta D$ and an endo-permutation source.
This source is constructed as described in Remark~\ref{rem:ProofFong2ndRed}: it can be identified with the cap of the endo-permutation $kD$-module $\Res^{ND}_D(V)$, say $V_0$, and we obtain that $V_0\cong k$ by Lemma~\ref{lem:sourcesSimpleDihedral}(b) as it is a source of a simple module for the $2$-nilpotent (hence soluble) group $ND$ and $D$ is dihedral. This proves that the above Morita equivalence is a splendid Morita equivalence.  
The last claim is immediate from Theorem~\ref{thm:Fong2ndReduction}(c). 
\end{proof}



\noindent Next, as a consequence of \cref{lem:coveringOddIndex}, we can describe all maximal defect $2$-blocks with dihedral defect groups of the finite groups $G$ with $O_{2^{\prime}}(G)=1$ up to splendid Morita equivalence. 

\begin{prop}\label{prop:ClassificBlocksDihedralDefectGroupsMaxDefectO_2prime_is_1}%
Let $G$ be a finite group with dihedral Sylow $2$-subgroups of order $2^n$ $(n\geq 3)$. Write $q = r^f$ for a power of a prime number $r$ with $f$ a positive
integer. If $O_{2'}(G)=1$, then any block $\bB\in\Bl^{\mathrm{max}}(kG)$  is splendidly Morita equivalent to $\bB_0(kG)$. 
Moreover, the following assertions hold.
    \begin{enumerate}[label={\rm(\alph*)}]
        \item If $G=D_{2^n}$ with $n\geq 3$, then $\Bl^{\mathrm{max}}(kG)=\{\bB_0(kG)\}$.
        \item If $n=3$ and $G=\fA_7$, then $\Bl^{\mathrm{max}}(kG)=\{\bB_0(kG)\}$.
        \item If $G=\PSL_2(q)\rtimes C_d$ with $(q-1)_2=2^n$ or $(q+1)_2=2^n$ and $d$ is an odd integer such that $d\,|\,f$, 
        then $|\Bl^{\mathrm{max}}(kG)|=d$ and any  $\bB\in\Bl^{\mathrm{max}}(kG)$ is splendidly Morita equivalent to $\bB_0(k\!\PSL_2(q))$.
        \item If $G=\PGL_2(q)\rtimes C_d$ with $(q-1)_2=2^{n-1}$ or $(q+1)_2=2^{n-1}$ and $d$ is an odd integer such that $d\,|\,f$, 
        then $|\Bl^{\mathrm{max}}(kG)|=d$ and any  $\bB\in\Bl^{\mathrm{max}}(kG)$ is splendidly Morita equivalent to $\bB_0(k\!\PGL_2(q))$.
    \end{enumerate}
 In cases {\rm(c)} and {\rm(d)}, any choice of q satisfying the indicated condition yields a representative of the corresponding splendid Morita equivalence class.      
\end{prop}

\begin{proof}
First, it follows from Subsection~\ref{ssec:dihedralSylow2} that $G$ is one of the groups listed in (a) to (d). 
Let $D\in\Syl_2(G)$. 
The number of blocks of maximal defect is easily calculated as follows. 
First, as $D$ is dihedral of order at least $8$, its automorphism group is a $2$-group, so
\[
N_G(D)= DC_G(D) =D\times O_{2^\prime}(N_G(D))   
\text{ and } 
C_G(D)=Z(D)\times O_{2^\prime}({N_G(D)})\,.
\]
Then  \cite[Lemma 2.1]{NavarroTiep12} shows (as a corollary of Brauer's 1st Main Theorem)  that  {$|\Bl^{\mathrm{max}}(kG)|$} is the number of $N_G(D)$-orbits on the conjugacy classes of elements in $O_{2^\prime}({N_G(D)})$.
\par
Now, Case (a) is clear. Indeed, as $G=D$ is a $2$-group, we have $kG=\bB_0(kG)$. Case~(b) is also clear as $D$ is self-normalising in $G=\fA_7$, and so the principal block is the unique block of maximal defect by the above.\par
Thus, we can assume that $G$ is one of the groups in (c) or (d).  Set $G_0:=\PSL_2(q)$ in Case~(c) and $G_0:=\PGL_2(q)$ in Case (d). Then $D\in\Syl_2(G_0)$ and  we have $N_{G_0}(D)=D$. See~\cite[Theorem~1.4.3(d)]{BonnafeBookSL2q} for $\PSL_2(q)$. 
For $G_0=\PGL_2(q)$ this easily follows from the fact that we may choose $D$ to be the Sylow $2$-subgroup of the normalizer of the split (resp. non-split) torus $T$ when $q\equiv 1\pmod{4}$ (resp. $q\equiv 3\pmod{4}$) and $N_{G_0}(T)\cong D_{2(q-1)}$ (resp. $N_{G_0}(T)\cong D_{2(q+1)}$). Any element of $N_{G_0}(D)$ must then normalize the unique cyclic subgroup of $D$ of index $2$ and it follows that $N_{G_0}(D)\leq N_{G_0}(T)$ and $N_{G_0}(D)=D$.
It follows that the principal block is the unique block of maximal defect of $kG_0$. If $d>1$, then $N_G(D)=D\times C_d$ and we obtain that $|\Bl^{\mathrm{max}}(kG)|=d$ by the above counting argument. Moreover, any block $\bB\in\Bl^{\mathrm{max}}(kG)$ 
covers a block $\bb\in\Bl^{\mathrm{max}}(kG_0)=\{\bB_0(kG_0)\}$ (as a defect group of $\bb$ is $G$-conjugate to $D\cap G_0=D$). By Lemma~\ref{lem:coveringOddIndex}
any such block is splendidly Morita equivalent to $\bB_0(kG_0)$. In particular, all the elements of $\Bl^{\mathrm{max}}(kG)$ are splendidly Morita equivalent to $\bB_0(kG)$.
Finally,  in cases {\rm(c)} and {\rm(d)}, the fact that the corresponding splendid Morita equivalence class does not depend on $q$ is given by \cite[Theorem~1.1]{KoLa19}.
The claim follows.
\end{proof}




